\scanpage{77}
\mspage{83}

\deleted{Cela doit signifier aussi que la \unsure{convolution|construction|considération|conclusion} de} Signalons

\uline{Proposition} 4. Soit $f:X\longrightarrow Y$ \uline{propre}. Pour que $f$ soit une $\mathbb D$-fibration, il f. et s. que pour \deleted{tout} $y_1\longrightarrow y_0$ de $Y$, le foncteur $\deleted{\unsure{B}}\ X_{/y_1}\unsure{-|\longrightarrow}X_{/y_0}$ induit soit une $\mathbb D$-équiv. \inserted{(i.e que $f$ soit une $\mathbb D$-fibr. \uline{coh.})}. Dualement, si $f$ est $\mathbb D$-lisse, pour que ce soit une $\mathbb D$-fibration, il \inserted{f. et} suffit que $\forall\, y_1\longrightarrow y_0$ de $Y$, $ {}_{y_0\backslash}X\longrightarrow {}_{y_1\backslash}X$ soit une $\mathbb D$-équivalence.

\margin{\uline{Cor} \deleted{\unsure{$f$}} Supposons $f$ Cat-cofibrant (donc propre). Pour que $f$ soit une $\mathbb D$-fibration, il f. et s. que les foncteurs $X_{y_0}\longrightarrow X_{y_1}$ (pour $y_0\longrightarrow y_1$ de $Y$) soient des $\mathbb D$-équivalences. Dualement pour Cat-fibration (donc foncteurs lisses)}

On pourrait appeler $\mathbb D$-fibrations \uline{parfaites} les foncteurs à la fois propres et lisses qui sont des \inserted{$\mathbb D$-}fibr. coh. (ou ce qui revient au $=$, des $\mathbb D$-fibr. \unsure{droites|dro.}).

\deleted{Définition} \uline{Proposition} 5. Soit $f:X\longrightarrow Y$ \unsure{de|dans} $\mathrm{Diag}$. Conditions équivalentes

a) \deleted{\unsure{$\forall$}} $f$ est une $\mathbb D$ équiv., et le reste par tt changement de base

b) $f$ est une $\mathbb D$ fibr. à fibres \inserted{$\mathbb D$-}contractiles

\uline{Dém} $a\Longrightarrow b$ Que \deleted{\unsure{$f$}} a) implique que $f$ soit une $\mathbb D$-fibr. est immédiat sur la définition de celle-ci (condition a) de la prop. 2), et que les fibres sont contractiles est tautologique.

$b)\Longrightarrow a)$ Il suffit de montrer que b) implique que $f$ est une $\mathbb D$ équivalence, Mais c'est un cas particulier de Prop. 3 b) (en y prenant $X'=Y$).
% archive sheet number: 76
